\section{Institution-Specific Constraints}
\label{sec:constraints}

\subsection{Inherited constraint mapping}

The Semester~VI phase identified the institution's scheduling rules and established that each is
expressible within UniTime's constraint model. Table~\ref{tab:inherited} reproduces that mapping,
which remains valid and is the starting point for this section.

\begin{table}[H]
\centering
\caption{Constraint mapping established in Semester VI.}
\label{tab:inherited}
\footnotesize
\begin{tabular}{@{}L{5.0cm}L{9.9cm}@{}}
\toprule
\textbf{Constraint} & \textbf{Mechanism in UniTime} \\
\midrule
Distance between rooms & Distance check integrated into the room detail configuration. \\
Preferred labs and lectures & Subparts within an instructional offering, each with its own time pattern. \\
Evenly distributed working hours & Instructor preference settings; the solver favours consecutiveness. \\
Minimisation of schedule holes & The \texttt{StudentMinimizeScheduleHoles} optimisation criterion. \\
Consecutive lectures for a course & Distribution preferences favouring adjacent slots on the same day. \\
Maximum slots per course per day & Distribution preferences imposing a per-day limit. \\
Instructor load management & Distribution preferences limiting back-to-back classes. \\
Departmental cohort conflicts & Distribution preferences preventing overlap for a shared cohort. \\
Specialised laboratory allocation & Room groups encoding hardware and equipment requirements. \\
Common university hours & Reserved slots excluded from the time-pattern configuration. \\
\bottomrule
\end{tabular}
\end{table}

\subsection{From mapping to implementation}

A mapping records that a rule \emph{can} be expressed; it does not schedule anything. The work of
this phase was to generate the corresponding data for every class in the dataset, so that the rules
are reproduced automatically each time the pipeline runs rather than re-entered by hand after each
import. Table~\ref{tab:implemented} lists what is actually emitted, and where.

\begin{table}[H]
\centering
\caption{Constraints as implemented in the generated data.}
\label{tab:implemented}
\footnotesize
\begin{tabular}{@{}L{3.6cm}L{5.0cm}L{6.2cm}@{}}
\toprule
\textbf{Rule} & \textbf{Where it is produced} & \textbf{Form} \\
\midrule
Meeting patterns & \texttt{gen\_preferences.py} & a \texttt{<timePref>} naming one of the ten
  session time patterns, on every subpart and class \\
Date pattern & \texttt{gen\_preferences.py} & \texttt{Full Term} date preference at required level \\
Minor-degree block & \texttt{FIXED\_ENTRIES\_MAP} & required Monday/Tuesday 16:30--18:30 window \\
Open-elective block & \texttt{FIXED\_ENTRIES\_MAP} & required Monday/Tuesday 09:30--10:30 window \\
Protection of both blocks & \texttt{FIXED\_ENTRIES\_MAP} & prohibited windows on all other courses \\
Room preferences & \texttt{gen\_preferences.py} & \texttt{<roomPref>} lists derived from Taasika
  room associations, never pre-assignments \\
Specialised laboratories & \texttt{gen\_buildings.py} & building inference plus a separate pool for
  unavailable laboratories \\
Lunch breaks & \texttt{setup\_lunch\_breaks.py} & solver criteria, applied server-side \\
Division and batch structure & \texttt{gen\_academic.py}, \texttt{gen\_students.py} & 242 student
  groups; membership derived from enrolment \\
Weekend exclusion & \texttt{gen\_buildings.py} & room sharing marks weekends unavailable \\
Travel times & \texttt{gen\_buildings.py} & 5 minutes between buildings, 3 between the two CSE
  buildings \\
\bottomrule
\end{tabular}
\end{table}

\subsection{The modular time-block mechanism}

The institution reserves fixed windows in the week for course categories that are taught across
departments. Minor-degree courses meet on Monday and Tuesday evenings; open electives meet on
Monday and Tuesday mornings. These are not preferences but hard institutional commitments, and
they have a second half that is easy to overlook: if a window is reserved for one category, every
\emph{other} course must be kept out of it, or the reservation achieves nothing.

This was implemented as a small rule engine inside \texttt{gen\_preferences.py}. A rule table,
\texttt{FIXED\_ENTRIES\_MAP}, is consulted for every subpart; the first rule whose course pattern
matches, and whose exclusion pattern does not, becomes the active rule for that subpart.
Table~\ref{tab:fixedentries} lists the four rules currently in force.

\begin{table}[H]
\centering
\caption{The fixed time-block rules, in evaluation order.}
\label{tab:fixedentries}
\footnotesize
\begin{tabular}{@{}L{3.1cm}L{5.2cm}L{1.9cm}L{2.3cm}L{1.9cm}@{}}
\toprule
\textbf{Rule} & \textbf{Courses matched} & \textbf{Days} & \textbf{Window} & \textbf{Level} \\
\midrule
\texttt{MDM\_Required}        & \texttt{MDM*}                                   & Mon, Tue & 16:30--18:30 & Required \\
\texttt{OE\_Required}         & \texttt{OE*}                                    & Mon, Tue & 09:30--10:30 & Required \\
\texttt{Non\_MDM\_Prohibited} & everything except \texttt{MDM*}, \texttt{OE*}   & Mon, Tue & 16:30--18:30 & Prohibited \\
\texttt{SY\_Non\_OE\_Prohibited} & CO, CoI, DTL, DTL-Lab, Eco, OOPD, OOPD-Lab, TOC, TOC-Tut & Mon, Tue & 09:30--10:30 & Prohibited \\
\bottomrule
\end{tabular}
\end{table}

The two levels do different jobs. A required rule marks the one slot that matches the day set and
fits the window; a prohibited rule marks every candidate whose days and interval overlap the
window. The first pins a course into its block, the second clears the block for everyone else.
The second-year rule lists courses by name rather than by wildcard, because only that year's core
courses share the open-elective morning, and the list says so explicitly.

\subsection{Meeting patterns}

Each subpart's weekly contact minutes, read from the Taasika snapshot, determine its meeting
pattern. Lectures of 240 minutes per week become \texttt{2 x 120}, 180 become \texttt{3 x 60} and
120 become \texttt{2 x 60}; laboratories and recitations are single meetings, so 480 minutes becomes
\texttt{1 x 480 Project}, 180 becomes \texttt{1 x 180} and 120 becomes \texttt{1 x 120}. Anything
that does not match a declared pattern falls back to \texttt{Exact Time}, which
\texttt{verify\_time\_patterns.py} reports as an audit finding rather than passing silently. The
migration of lectures to the $2\times120$ pattern was a deliberate change made during this phase to
reflect how the department actually teaches.

\subsection{Lunch breaks}

The Semester~VI report proposed implementing a universal lunch break by excluding slots from the
assignable time patterns. That approach was rejected in implementation because it is too rigid: it
removes the same hour from every day for every class, whereas the institution requires only that
each student and each instructor have \emph{some} free hour within a window.

The implemented mechanism is a dynamic one-hour break between 11:30 and 15:30, applied to students
and to instructors, and installed as solver criteria by the Jython script
\texttt{scripts/unitime\_scripts/setup\_lunch\_breaks.py}, which is uploaded and run inside UniTime
through \textbf{Administration\,$\rightarrow$\,Utilities\,$\rightarrow$\,Scripts}. Rooms remain
available for the whole weekday. Because this is a soft criterion rather than a blocked slot, the
solver may violate it, and the evaluator of Section~\ref{sec:evaluation} therefore checks it
explicitly: the final solution leaves 63 students and one instructor without a qualifying break.

\subsection{Student groups and what they do not do}

The dataset declares 242 student groups, one for each Taasika class and each batch. Since the
pipeline's most recent change, every student is also made a member of the groups corresponding to
the classes they are enrolled in: the group code is taken from the class's schedule note, so a
student enrolled in the second-year division-one lecture is a member of \texttt{SY-Div1}, and a
student enrolled in a particular laboratory batch is a member of that batch's group.

Membership makes the cohort visible and is a prerequisite for a future group reservation. It does
not constrain the timetable. A student's groups affect course sectioning in one place only, a
distance computation used for students who are not yet enrolled, and the final sectioning pass does
not read them. With student switching off, adding the memberships changed no placement and no
conflict count. What actually pins a student to a section is file 16 together with the three
settings of Section~\ref{sec:setup}.

\subsection{Shared laboratory sessions}

A laboratory supervised by two instructors, in the same room at the same time for the same batch,
is represented through UniTime's support for multiple lead instructors on a single class. Courses
with several lecture or laboratory sections taught by different instructors are modelled as
separate classes under a common subpart, which is the same mechanism the pipeline uses to emit one
lecture class per Taasika class record and one laboratory class per batch record.
